\section{Relation to Optimized Regression Discontinuity Designs}\label{sec:appendix:iw}

This appendix discusses the relationship between the segmented estimator and the optimized RD estimator of \textcite{imbens_optimized_2019}, focusing on why the choice of smoothness class determines whether segmentation can improve inference.


\subsection*{The Imbens--Wager Smoothness Class}

\textcite{imbens_optimized_2019} assume that the conditional response functions satisfy a bounded second derivative condition.
In the multivariate case with $X_i \in \mathbb{R}^k$, this takes the form $\|\nabla^2 \mu_w(x)\| \leq B$ for $w \in \{0,1\}$, where $\|\cdot\|$ denotes the operator norm (largest absolute eigenvalue of the Hessian).
This is an isotropic smoothness assumption: it constrains curvature in all directions jointly through a single constant $B$.

The present paper works with a first-order (Lipschitz) condition rather than a second-order condition, motivated by the use of the local constant estimator.
The isotropic Lipschitz analog is $\mathcal{F}_1^{\mathrm{iso}}(M)$ defined in \cref{sec:appendix:isotropic}.
The arguments below apply at both orders.


\subsection*{Why Segmentation Cannot Improve on Pooling under the Isotropic Class}

Under the isotropic class, the gradient of $f(x,z)$ is constrained to lie in a disk:
\begin{equation*}
	\left(\frac{\partial f}{\partial x}\right)^2 + \left(\frac{\partial f}{\partial z}\right)^2 \leq M^2.
\end{equation*}
The curvature budget is \emph{fungible} across directions.
Segmentation controls the $z$-channel of the bias by restricting comparisons to narrow segments, making the estimator insensitive to variation in $z$.
However, once the $z$-channel is shut down, the worst-case function simply reallocates its entire curvature budget to the $x$-direction: setting $\partial f / \partial z = 0$ and $\partial f / \partial x = M$.
The resulting worst-case bias is $h \cdot M \cdot B_{1,0}$, which is the same as the worst-case bias of the pooled estimator --- because the pooled estimator's worst case also concentrates curvature in whichever direction maximizes the bias.

In other words, under the isotropic class, segmentation does not reduce worst-case bias: any curvature removed from the compositional channel reappears in the smoothing channel.


\subsection*{Why Anisotropic Smoothness is Necessary}

Under the anisotropic class $\mathcal{F}_1^{\mathrm{aniso}}(M_x, M_z)$, the partial derivative bounds are \emph{separate}:
\begin{equation*}
	\abs*{\frac{\partial f}{\partial x}} \leq M_x, \qquad \abs*{\frac{\partial f}{\partial z}} \leq M_z.
\end{equation*}
The adversary cannot redirect the $M_z$ budget into the $x$-direction.
When segmentation neutralizes the $z$-channel, the worst-case bias falls from a quantity depending on both $M_x$ and $M_z$ (under pooling) to one depending primarily on $M_x$ (under segmentation).
The $M_z$ budget is stranded --- it can only generate within-segment compositional variation proportional to $M_z \cdot \delta$, which is small for narrow segments.

This is precisely why the anisotropic class is the natural framework for studying the benefits of covariate adjustment via segmentation: it is the weakest class under which controlling for the covariate direction yields a strict improvement in worst-case bias.


\subsection*{Minimax Optimal Estimation}

\begin{definition}[Minimax Optimal Linear Estimator]\label{def:minimax}
	Let $\mathcal{L}$ denote the class of linear estimators $\hat\tau = \sum_{i=1}^n w_i Y_i$.
	For a linear estimator with fixed weights, the worst-case mean squared error over a function class $\mathcal{F}$ is
	\begin{equation}\label{eq:minimax_mse}
		\sup_{f \in \mathcal{F}} \mathrm{MSE}(\hat\tau, f) = \wcb^2(\hat\tau, \mathcal{F}) + \Var(\hat\tau),
	\end{equation}
	where $\wcb(\hat\tau, \mathcal{F}) = \sup_{f \in \mathcal{F}} \abs{\E[\hat\tau] - \tau(f)}$ is the worst-case bias.
	The variance does not depend on $f$ because $\hat\tau$ is linear with fixed weights.
	The minimax optimal linear estimator over $\mathcal{F}$ is
	\begin{equation}\label{eq:minimax_estimator}
		\hat\tau^*_{\mathcal{F}} = \argmin_{\hat\tau \in \mathcal{L}} \left[\wcb^2(\hat\tau, \mathcal{F}) + \Var(\hat\tau)\right].
	\end{equation}
\end{definition}

\textcite{imbens_optimized_2019} characterize $\hat\tau^*$ for the isotropic class (with a second-order smoothness condition).
Their estimator is minimax optimal among all linear estimators under that class.


\subsection*{The Anisotropic Class as the Natural Framework for Covariate Adjustment}

The key question is not which smoothness class is ``correct,'' but which class captures the structure relevant for studying covariate adjustment.

\paragraph{Controlling $z$ should not worsen $x$.}
When an estimator controls for the $z$-direction --- restricting comparisons to observations with similar $z$-values --- one would not expect the regression function to become harder to estimate in the $x$-direction.
The smoothness of $m(x,z)$ in $x$ is a property of the data-generating process, not something that responds to the researcher's estimation strategy.
The anisotropic class respects this: $\abs{\partial f / \partial x} \leq M_x$ holds independently of what happens in the $z$-direction.
The isotropic class violates it: controlling for $z$ frees up curvature budget that the worst-case function reallocates to $x$, making the $x$-direction effectively harder.

\paragraph{Minimax implications.}
The nesting $\mathcal{F}_1^{\mathrm{aniso}}(M, M) \subsetneq \mathcal{F}_1^{\mathrm{iso}}(M\sqrt{2})$ (\cref{sec:appendix:isotropic}) implies that the minimax risk under the anisotropic class is weakly lower:
\begin{equation}\label{eq:minimax_comparison}
	\min_{\hat\tau \in \mathcal{L}} \sup_{f \in \mathcal{F}_{\mathrm{aniso}}(M, M)} \mathrm{MSE}(\hat\tau, f) \;\leq\; \min_{\hat\tau \in \mathcal{L}} \sup_{f \in \mathcal{F}_{\mathrm{iso}}(M\sqrt{2})} \mathrm{MSE}(\hat\tau, f).
\end{equation}
This inequality is mechanical (monotonicity of minimax risk in the function class).
What is not mechanical is \emph{why} the inequality should be strict: the anisotropic class has more structure --- two separate directional bounds --- that the minimax estimator can exploit.
Under $\mathcal{F}_{\mathrm{iso}}(M\sqrt{2})$, segmentation is pointless because the curvature budget is fungible; the minimax estimator ignores $z$.
Under $\mathcal{F}_{\mathrm{aniso}}(M, M)$, the separate bounds prevent curvature reallocation, so the minimax estimator can incorporate $z$-information via segmentation or a similar mechanism, achieving a strictly lower worst-case MSE whenever $M > 0$.

\paragraph{Practical consequence.}
The segmented estimator is not minimax optimal under the anisotropic class, but it reduces worst-case bias relative to pooling (\cref{sec:bias_comparison}) and is simple to implement --- it corresponds to the segment fixed effects already common in applied work.
Characterizing the exact minimax estimator under the anisotropic class is an open problem; the present paper establishes that the anisotropic class is the appropriate framework and that segmentation yields a strict improvement within it.


\subsection*{Balancing Constraints and the Separability Argument}

In Section~V.B, \textcite{imbens_optimized_2019} propose incorporating auxiliary covariates $Z_i \in \mathbb{R}^p$ by adding \emph{balancing constraints} $\sum_{i=1}^n \hat\gamma_i Z_{ij} = 0$ for $j = 1, \ldots, p$ to the weight optimization.
These constraints force the estimator weights to be balanced in the covariate directions, and are motivated as a robustness device: even if the smoothness class does not model $Z$, enforcing balance ``may improve robustness'' to hidden biases and can reduce variance when $Y_i$ depends additively on $Z_i$.

The separability framework developed in this paper reveals why such external constraints are necessary under Imbens--Wager's setup --- and why they cannot improve honest inference.

\paragraph{Why the constraints are necessary.}
The Imbens--Wager smoothness class constrains $\mu_w(x)$ as a function of the running variable $X$ only.
Covariates $Z$ live entirely outside the function class: the minimax machinery is blind to $z$-variation.
There is no mechanism within the smoothness class to encode the covariate's role, so covariate information must be injected externally through weight constraints.
By the separability trilemma (\cref{cor:iso_inflates}), extending the smoothness class to include $z$ isotropically --- e.g., $\norm{\nabla^2 m(x,z)}_{\mathrm{op}} \leq B$ --- would not resolve this.
The operator-norm class couples the primitive bounds: it constrains the curvature in $x$ and $z$ jointly through a single constant $B$.
To remain coupling-agnostic with respect to the primitive entry-wise bounds $(B_{xx}, B_{xz}, B_{zz})$, the operator-norm ball must inflate to radius $nB$ (\cref{app:modulus:hessian_norms}), strictly exceeding the primitive bounds.
The resulting function class is strictly larger than the separable class --- the curvature budget is fungible across directions, and the worst-case function concentrates the entire budget into whichever direction maximizes the bias, rendering covariate adjustment pointless.

\paragraph{Why the constraints cannot improve honest inference.}
Under the isotropic class, the unconstrained minimax estimator already ignores $z$: the worst-case function sets $\partial f / \partial z = 0$ and concentrates the entire curvature budget into $x$.
Adding balancing constraints restricts the class of estimators (from $\mathcal{L}$ to $\mathcal{L}_{\mathrm{bal}} \subsetneq \mathcal{L}$) while leaving the function class unchanged.
Since the minimax risk is a minimum over estimators,
\begin{equation*}
	\min_{\hat\tau \in \mathcal{L}_{\mathrm{bal}}} \sup_{f \in \mathcal{F}_{\mathrm{iso}}} \mathrm{MSE}(\hat\tau, f) \;\geq\; \min_{\hat\tau \in \mathcal{L}} \sup_{f \in \mathcal{F}_{\mathrm{iso}}} \mathrm{MSE}(\hat\tau, f).
\end{equation*}
The balancing constraints can only weakly \emph{worsen} worst-case performance.
They may reduce empirical bias when the true regression function is better behaved than the worst case, but the worst-case bias bound --- which governs honest confidence intervals --- is unchanged or inflated.

\paragraph{The separable class resolves both problems.}
The separable class $\mathcal{F}(\mathcal{C}_{\mathrm{rect}})$ models the covariate direction directly through separate bounds $(M_x, M_z)$.
The $z$-budget is stranded: the adversary cannot reallocate $M_z$ into the $x$-direction.
The minimax estimator over $\mathcal{F}(\mathcal{C}_{\mathrm{rect}})$ automatically incorporates covariate information --- it downweights cross-$z$ comparisons because doing so reduces worst-case bias, not as an external requirement.
Covariate adjustment is a consequence of minimax optimality under the separable class, not an ad hoc correction.

The practical implication is that a researcher who has primitive bounds $(M_x, M_z)$ and works with the separable class obtains:
\begin{enumerate}
	\item A strictly smaller function class than the coupling-agnostic isotropic class (\cref{cor:iso_inflates}), and hence a strictly lower minimax risk.
	\item An estimator that uses covariate information through the minimax machinery, yielding honest confidence intervals that are genuinely shorter --- not merely empirically shorter.
	\item No need for external balancing constraints: the structure that Imbens--Wager inject through weight restrictions is already encoded in the function class.
\end{enumerate}